\chapter{Logging}
\label{CH:LOG}
%         is it process or kernel thread?
% 
% 	be sure to say buffer, not block 
% 
% 	Mount

%% 
%%  -------------------------------------------
%% 

One of the most interesting problems in file system design is crash
recovery. The problem arises because many file-system operations
involve multiple writes to the disk, and a crash after a subset of the
writes may leave the on-disk file system in an inconsistent state. For
example, suppose a crash occurs during file truncation.
Truncation involves writing (at least) two disk blocks,
one containing the i-node and one containing part of the free-block bitmap:
the i-node's length must be set to zero and the block
numbers in the i-node must be cleared,
and the bitmap bits for the file's blocks must
be set to zero (marking them free).
A crash after one of the writes, but before the kernel has
a chance to perform the second write, will leave the
on-disk file system in an incorrect state.
If only the bitmap write occurs, the result after a crash
and reboot will be an i-node with non-zero length which
refers to blocks that are marked free;
if only the i-node write occurs, the result will be
blocks that are not used by any file but are also not marked free.

The latter is relatively benign, but an inode that refers to a freed
block is likely to cause serious problems after a reboot.  After reboot, the
kernel might allocate that block to another file, and now we have two different
files pointing unintentionally to the same block.  If xv6 supported
multiple users, this situation could be a security problem, since the
old file's owner would be able to read and write blocks in the
new file, owned by a different user.

Xv6 solves the problem of crashes during file-system operations with a
simple form of logging. An xv6 system call does not directly write
the on-disk file system data structures. Instead, it places a
description of all the disk writes it wishes to make in a 
\indextext{log} 
on the disk. Once the system call has logged all of its writes, it writes a
special 
\indextext{commit}
record to the disk indicating that the log contains
a complete operation. At that point the system call copies the writes
to the on-disk file system data structures. After those writes have
completed, the system call erases the log on disk.

If the system should crash and reboot, the file-system code recovers
from the crash as follows, before running any processes. If the log is
marked as containing a complete operation, then the recovery code
copies the writes to where they belong in the on-disk file system. If
the log is not marked as containing a complete operation, the recovery
code ignores the log.  The recovery code finishes by erasing
the log.

Why does xv6's log solve the problem of crashes during file system
operations? If the crash occurs before the operation commits, then the
log on disk will not be marked as complete, the recovery code will
ignore it, and the state of the disk will be as if the operation had
not even started. If the crash occurs after the operation commits,
then recovery will replay all of the operation's writes, perhaps
repeating them if the operation had started to write them to the
on-disk data structure. In either case, the log makes operations
atomic with respect to crashes: after recovery, either all of the
operation's writes appear on the disk, or none of them appear.

%% 
%% 
%% 

\section{Buffer cache layer}
\label{s:bcache}

The logging layer (and the file system) rely on the buffer cache to
read and write disk blocks (see Figure~\ref{fig:fslayer}).
The buffer cache has two jobs: (1) synchronize access to disk blocks to ensure
that only one copy of a block is in memory and that only one kernel thread at a time
uses that copy; (2) cache popular blocks so that they don't need to be re-read from
the slow disk. The code is in
\lstinline{bio.c}.

The main interface exported by the buffer cache consists of
\indexcode{bread}
and
\indexcode{bwrite};
the former obtains a
\indextext{buf}
containing a copy of a block which can be read or modified in memory, and the
latter writes a modified buffer to the appropriate block on the disk.
A kernel thread must release a buffer by calling
\indexcode{brelse}
when it is done with it.
The buffer cache uses a per-buffer sleep-lock to ensure
that only one thread at a time uses each buffer
(and thus each disk block);
\lstinline{bread}
returns a locked buffer, and
\lstinline{brelse}
releases the lock.

The buffer cache has a fixed number of buffers to hold disk blocks,
which means that if the file system asks for a block that is not
already in the cache, the buffer cache must recycle a buffer currently
holding some other block. The buffer cache recycles the
least recently used buffer for the new block. The assumption is that
the least recently used buffer is the one least likely to be used
again soon.

There is a subtle interaction between the logging layer and the buffer
cache:  the buffer cache cannot recycle a buffer that is in use by a
transaction until the transaction commits (see
Section~\ref{s:code-logging}).
To avoid recycling too early, the logging layer uses
the buffer-cache functions \indexcode{bpin} and \indexcode{bunpin} to
pin and unpin a buffer in the buffer cache.

%% 
%%  -------------------------------------------
%% 
\section{Code: Buffer cache}

The buffer cache is a doubly-linked list of buffers.
The function
\indexcode{binit},
called by
\indexcode{main}
\lineref{kernel/main.c:/binit/},
initializes the list with the
\indexcode{NBUF}
buffers in the static array
\lstinline{buf}
\linerefs{kernel/bio.c:/Create.linked.list/,/^..}/}.
All other access to the buffer cache refer to the linked list via
\indexcode{bcache.head},
not the
\lstinline{buf}
array.

A buffer has two state fields associated with it.
The field
\indexcode{valid}
indicates that the buffer contains a copy of the block.
The field \indexcode{disk}
indicates that the buffer content has been handed to
the disk, which may change the buffer (e.g., write
data from the disk into \lstinline{data}).

\lstinline{bread}
\lineref{kernel/bio.c:/^bread/}
calls
\indexcode{bget}
to get a buffer for the given block
\lineref{kernel/bio.c:/b.=.bget/}.
If the buffer needs to be read from disk,
\lstinline{bread}
calls
\indexcode{virtio_disk_rw}
to do that before returning the buffer.

\lstinline{bget}
\lineref{kernel/bio.c:/^bget/}
scans the buffer list for a buffer with the given device and block numbers
\linerefs{kernel/bio.c:/Is.the.block.already/,/^..}/}.
If there is such a buffer,
\indexcode{bget}
acquires the sleep-lock for the buffer.
\lstinline{bget}
then returns the locked buffer.

If there is no cached buffer for the given block,
\indexcode{bget}
must make one, possibly reusing a buffer that held
a different block.
It scans the buffer list a second time, looking for a buffer
that is not in use (\lstinline{b->refcnt = 0});
any such buffer can be used.
\lstinline{bget}
edits the buffer metadata to record the new device and block number
and acquires its sleep-lock.
Note that the assignment
\lstinline{b->valid = 0}
ensures that
\lstinline{bread}
will read the block data from disk
rather than incorrectly using the buffer's previous contents.

It is important that there is at most one cached buffer per
disk block, to ensure that readers see writes, and because the
file system uses locks on buffers for synchronization.
\lstinline{bget}
ensures this invariant by holding the
\lstinline{bcache.lock}
continuously from the first loop's check of whether the
block is cached through the second loop's declaration that
the block is now cached (by setting
\lstinline{dev},
\lstinline{blockno},
and
\lstinline{refcnt}).
This causes the check for a block's presence and (if not
present) the designation of a buffer to hold the block to
be atomic.

It is safe for
\lstinline{bget}
to acquire the buffer's sleep-lock outside of the 
\lstinline{bcache.lock}
critical section,
since the non-zero
\lstinline{b->refcnt}
prevents the buffer from being re-used for a different disk block.
The sleep-lock protects reads
and writes of the block's buffered content, while the
\lstinline{bcache.lock}
protects information about which blocks are cached.

If all the buffers are busy, then too many processes are
simultaneously executing file system calls;
\lstinline{bget}
panics.
A more graceful response might be to sleep until a buffer became free,
though there would then be a possibility of deadlock.

Once
\indexcode{bread}
has read the disk (if needed) and returned the
buffer to its caller, the caller has
exclusive use of the buffer and can read or write the data bytes.
If the caller does modify the buffer, it must call
\indexcode{bwrite}
to write the changed data to disk before releasing the buffer.
\lstinline{bwrite}
\lineref{kernel/bio.c:/^bwrite/}
calls
\indexcode{virtio_disk_rw}
to talk to the disk hardware.

When the caller is done with a buffer,
it must call
\indexcode{brelse}
to release it. 
(The name
\lstinline{brelse},
a shortening of
b-release,
is cryptic but worth learning:
it originated in Unix and is used in BSD, Linux, and Solaris too.)
\lstinline{brelse}
\lineref{kernel/bio.c:/^brelse/}
releases the sleep-lock and
moves the buffer
to the front of the linked list
\linerefs{kernel/bio.c:/b->next->prev.=.b->prev/,/bcache.head.next.=.b/}.
Moving the buffer causes the
list to be ordered by how recently the buffers were used (meaning released):
the first buffer in the list is the most recently used,
and the last is the least recently used.
The two loops in
\lstinline{bget}
take advantage of this:
the scan for an existing buffer must process the entire list
in the worst case, but checking the most recently used buffers
first (starting at
\lstinline{bcache.head}
and following
\lstinline{next}
pointers) will reduce scan time when there is good locality of reference.
The scan to pick a buffer to reuse picks the least recently used
buffer by scanning backward
(following 
\lstinline{prev}
pointers).
%% 
%%  -------------------------------------------
%% 

\section{Log design}

The log resides at a known fixed location (see
Figure~\ref{fig:fslayout}), specified in the superblock.
It consists of a header block followed by a sequence
of updated block copies (``logged blocks'').
The header block contains an array of block
numbers, one for each of the logged blocks, and 
the count of log blocks.
The count in the header block on disk is either
zero, indicating that there is no transaction in the log,
or non-zero, indicating that the log contains a complete committed
transaction with the indicated number of logged blocks.
Xv6 writes the header
block when a transaction commits, but not before, and sets the
count to zero after copying the logged blocks to the file system.
Thus a crash midway through a transaction will result in a
count of zero in the log's header block; a crash after a commit
will result in a non-zero count.

Each system call's code indicates the start and end of the sequence of
writes that must be atomic with respect to crashes.
To allow concurrent execution of file-system operations
by different processes,
the logging system can accumulate the writes
of multiple system calls into one transaction.
Thus a single commit may involve the writes of multiple
complete system calls.
To avoid splitting a system call across transactions, the logging system
only commits when no file-system system calls are underway.

The idea of committing several transactions together is known as 
\indextext{group commit}.
Group commit reduces the number of disk operations
because it amortizes the fixed cost of a commit over multiple
operations.
Group commit also hands the disk system more concurrent writes
at the same time, perhaps allowing the disk to write
them all during a single disk rotation.
Xv6's virtio driver doesn't support this kind of
\indextext{batching},
but xv6's file system design allows for it.

Xv6 dedicates a fixed amount of space on the disk to hold the log.
The total number of blocks written by the system calls in a
transaction must fit in that space.
This has two consequences.
No single system call
can be allowed to write more distinct blocks than there is space
in the log. This is not a problem for most system calls, but two
of them can potentially write many blocks: 
\indexcode{write}
and
\indexcode{unlink}.
A large file write may write many data blocks and many bitmap blocks
as well as an inode block; unlinking a large file might write many
bitmap blocks and an inode.
Xv6's write system call breaks up large writes into multiple smaller
writes that fit in the log,
and 
\lstinline{unlink}
doesn't cause problems because in practice the xv6 file system uses
only one bitmap block.
The other consequence of limited log space
is that the logging system cannot allow a system call to start
unless it is certain that the system call's writes will
fit in the space remaining in the log.
%% 
%% 
%% 
\section{Code: logging}
\label{s:code-logging}

A typical use of the log in a system call looks like this:
\begin{lstlisting}[]
  begin_op();
  ...
  bp = bread(...);
  bp->data[...] = ...;
  log_write(bp);
  ...
  end_op();
\end{lstlisting}

\indexcode{begin_op}
\lineref{kernel/log.c:/^begin.op/}
waits until
the logging system is not currently committing, and until
there is enough unreserved log space to hold
the writes from this call.
\lstinline{log.outstanding}
counts the number of system calls that have reserved log
space; the total reserved space is 
\lstinline{log.outstanding}
times
\lstinline{MAXOPBLOCKS}.
Incrementing
\lstinline{log.outstanding}
both reserves space and prevents a commit
from occurring during this system call.
The code conservatively assumes that each system call might write up to
\lstinline{MAXOPBLOCKS}
distinct blocks.

\indexcode{log_write}
\lineref{kernel/log.c:/^log.write/}
acts as a proxy for 
\indexcode{bwrite}.
It records the block's block number in memory,
reserving it a slot in the log on disk,
and pins the buffer in the block cache
to prevent the block cache from evicting it.
The block must stay in the cache until committed:
until then, the cached copy is the only record
of the modification; it cannot be written to
its place on disk until after commit;
and other reads in the same transaction must
see the modifications.
\lstinline{log_write}
notices when a block is written multiple times during a single
transaction, and allocates that block the same slot in the log.
This optimization is often called
\indextext{absorption}.
It is common that, for example, the disk block containing inodes
of several files is written several times within a transaction.  By absorbing
several disk writes into one, the file system can save log space and
can achieve better performance because only one copy of the disk block must be
written to disk.

\indexcode{end_op}
\lineref{kernel/log.c:/^end.op/}
first decrements the count of outstanding system calls.
If the count is now zero, it commits the current
transaction by calling
\lstinline{commit().}
There are four stages in this process.
\lstinline{write_log()}
\lineref{kernel/log.c:/^write.log/}
copies each block modified in the transaction from the buffer
cache to its slot in the log on disk.
\lstinline{write_head()}
\lineref{kernel/log.c:/^write.head/}
writes the header block to disk: this is the
commit point, and a crash after the write will
result in recovery replaying the transaction's writes from the log.
\indexcode{install_trans}
\lineref{kernel/log.c:/^install_trans/}
reads each block from the log and writes it to the proper
place in the file system.
Finally
\lstinline{end_op}
writes the log header with a count of zero;
this has to happen before the next transaction starts writing
logged blocks, so that a crash doesn't result in recovery
using one transaction's header with the subsequent transaction's
logged blocks.

\indexcode{recover_from_log}
\lineref{kernel/log.c:/^recover_from_log/}
is called from 
\indexcode{initlog}
\lineref{kernel/log.c:/^initlog/},
which is called from \indexcode{fsinit}\lineref{kernel/fs.c:/^fsinit/} during boot before the first user process runs
\lineref{kernel/proc.c:/fsinit/}.
It reads the log header, and mimics the actions of
\lstinline{end_op}
if the header indicates that the log contains a committed transaction.

An example use of the log occurs in 
\indexcode{filewrite}
\lineref{kernel/file.c:/^filewrite/}.
The transaction looks like this:
\begin{lstlisting}[]
      begin_op();
      ilock(f->ip);
      r = writei(f->ip, ...);
      iunlock(f->ip);
      end_op();
\end{lstlisting}
This code is wrapped in a loop that breaks up large writes into individual
transactions of just a few blocks at a time, to avoid overflowing
the log.  The call to
\indexcode{writei}
writes many blocks as part of this
transaction: the file's inode, one or more bitmap blocks, and some data
blocks.

%% 
%% %% 
%%  -------------------------------------------
%% 
\section{Real world}

The buffer cache in a real-world operating system is significantly
more complex than xv6's, but it serves the same two purposes:
caching and synchronizing access to the disk.
Xv6's buffer cache uses a least recently used (LRU)
eviction policy; there are many more complex
policies that can be implemented, each good for some
workloads and not as good for others.
A more efficient LRU cache would eliminate the linked list,
instead using a hash table for lookups and a heap for LRU evictions.
Modern buffer caches are typically integrated with the
virtual memory system to support memory-mapped files.

Xv6's logging system is inefficient.
A commit cannot occur concurrently with file-system system calls.
The system logs entire blocks, even if
only a few bytes in a block are changed. It performs synchronous
log writes, a block at a time, each of which is likely to require an
entire disk rotation time. Real logging systems address all of these
problems.
-
Logging is not the only way to provide crash recovery. Early file systems
used a scavenger during reboot (for example, the UNIX
\indexcode{fsck}
program) to examine every file and directory and the block and inode
free lists, looking for and resolving inconsistencies. Scavenging can take
hours for large file systems, and there are situations where it is not
possible to resolve inconsistencies in a way that causes the original
system calls to be atomic. Recovery
from a log is much faster and causes system calls to be atomic
in the face of crashes.


%%  -------------------------------------------
%% 
\section{Exercises}

\begin{enumerate}

  \item \mfk{TODO: add logging exercises}

\end{enumerate}
